\section{Computational results}
\label{sec:results}

\subsection{Protocol, instances and what each quantity means}
\label{sec:res-setup}

The masters, the partitions, the volumetric benchmark and the compact
models are solved with Gurobi 12; the pricing problems and the
placeability tests with CP-SAT from OR-Tools 9.12.
Table~\ref{tab:protocol} declares, per stage, the five parameters that
can change a solver's answer: time limit, MIP gap, threads, per-model
memory cap and seed. Memory is part of the budget --- Gurobi holds its
tree in memory --- so the time at which a fixed cap is reached licenses
no inference about difficulty. Every row below is read from a recorded run, each one declaring its
own budget, thread count, memory cap and commit. A \emph{gap} is
always $(z-\mathrm{LB})/z$ against the method's own certified bound;
Table~\ref{tab:quantities} defines the rest.

Every row is one run, and one run is reproducible: the procedure is
deterministic given a build, a machine and the way its code is
executed. That was checked rather than assumed, and the check is worth
reporting in full. On $I\_16\_15\_2\_90\_k25$ in 3D, three runs
with the solver's module compiled in the running process returned the
same plan at the same cost to the cent, and two runs of the same
source from its precompiled form also agreed to the cent --- with each
\emph{other}, at a different plan: $55$ trucks at $79\,608.45$
against $59$ at $83\,645.86$. Nothing about the model or the
configuration differs between them; only how fast the machine
executes the same instructions.

That is what a wall-clock budget costs. Every phase of the procedure
draws on one, so which phases finish, how many construction attempts
fit, and which regions the tree opens all depend on execution speed,
and a plan is the method's answer under one schedule rather than a
function of its configuration. The effect is largest where a phase
ends near its deadline: on the largest ladder instance the two
execution paths differ by whether a plan is found at all. Every run
reported here was executed the same way, so the table is internally
comparable, and two consequences are taken seriously throughout: a
difference of one or two trucks between rows is not read as an effect,
and no claim rests on one.

Three settings that are not budgets also decide what the tables can
say.
Column generation is run to a certificate at the root and at the one
region that can refute the incumbent outright --- ``is there a plan
with one truck fewer'' --- and nowhere else, which is the policy every
run reported here used. Certifying every node instead is affordable on
the volumetric model and closes far more of the space there, so the
\emph{open} counts below are a property of that choice and not of the
method; it is not affordable under 3D loading, where one geometric
sweep costs minutes, and the two modes are reported under one policy
so that their columns mean the same thing. The second is that the branching hierarchy's
third level --- the deliveries --- is disabled in every run reported
here, so a full-depth leaf is covered by a certificate rather than
split again; Section~\ref{sec:branching} says what that costs and what
it does not. The third is that no seed is set by the implementation of
Table~\ref{tab:volumetric} and the size ladder; nothing reported rests
on one.

\begin{table}[!htbp]
\centering
\small
\caption{Declared budgets, one line per stage that produces a number
reported below. \emph{gap} is the relative MIP gap at which an integer
solve may stop and \emph{mem} the per-model soft cap. The two
placeability censuses differ in oracle as well as in limit: the
thirty-six instances were put to the constraint programme alone, the
size ladder to a constructive shelf witness with the constraint
programme only for refutations, which is why only the first leaves
loads undecided. Every line but the last is read from the run files
themselves; the volumetric runs of BP-3D-SCPTOP recorded no
configuration and theirs is recovered from their launchers, including
the fact that a thread count asked for on the default Gurobi
environment did not reach the models, which therefore used all twelve
cores of the machine.}
\label{tab:protocol}
\resizebox{\textwidth}{!}{%
\begin{tabular}{llllrrl}
\toprule
Stage & solver & time (s) & gap & threads & mem (GB) & seed \\
\midrule
BP-3D-SCPTOP, 36 instances & Gurobi + CP-SAT & $900$ & --- & 10 & 20 & 0 \\
BP-3D-SCPTOP, size ladder  & Gurobi + CP-SAT & $900$--$3\,600$ & --- & 10 & 20 & 0 \\
Compact SCPTOP (volumetric) & Gurobi & $600$--$3\,600$ & $0$ & 10 & 20 & 0 \\
Compact 3D-SCPTOP, $n_{\mathrm{LC}}=3$ & Gurobi & $1\,800$ & $0$ & 10 & 20 & 0 \\
Fleet-floor certification & Gurobi & to convergence & --- & 1 & 20 & 0 \\
Placeability census, 36 instances & CP-SAT & $300$ per load & --- & --- & --- & --- \\
Placeability census, size ladder & shelf + CP-SAT & $60$ per load & --- & --- & --- & --- \\
\midrule
Fleet tree, volumetric oracle & Gurobi & $1\,800^{\sim}$ & --- & $12^{\sim}$ & $4^{\sim}$ & --- \\
\bottomrule
\end{tabular}}
\\[2pt]
{\footnotesize $^{\sim}$ recovered after the fact, not recorded; on one
of those five runs the time limit is not recoverable either.}
\end{table}


\begin{table}[!htbp]
\centering
\small
\caption{What each reported quantity means. A sentinel recorded by a
run that obtained no bound is never printed, never averaged and never
rounded to a side.}
\label{tab:quantities}
\begin{tabular}{@{}lp{0.80\textwidth}@{}}
\toprule
$z$ & The cost \eqref{eq:MPObjective} of the plan returned, net of any
  artificial penalty; a plan is reported only if it carries no
  artificial slack. An upper bound on the optimum in every case. \\
trucks & $\sum_{g,t}n_{gt}$ for that plan. \\
LB & A lower bound on the optimum of the model the method solves. For
  the method it is the \emph{certified} bound of
  Section~\ref{sec:bounds}, printed only where the sweep behind it came
  back with every pair proved; where it did not, the row says
  \emph{unproved} and carries no gap. For a compact model it is the
  solver's dual bound at its stop. \\
gap & $(z-\mathrm{LB})/z$, against the method's own bound. An
  optimality gap only to the extent that bound is tight. \\
stop & How a solve ended: \textsc{optimal}, \textsc{time\_limit} or
  \textsc{mem\_limit}. A time or memory stop makes the row's time the
  budget rather than a solve time. \\
nodes / open & Tree nodes opened; regions left unexpanded at the
  deadline. A run with open regions has not searched the space,
  whatever its gap --- and a run with none has not refuted it either,
  a region closed with a bound below the incumbent being bounded and
  not excluded. \\
verified & Dispatched trucks re-derived from the column's own fill and
  found legal by the independent routine of
  Section~\ref{sec:res-setup}, out of the number dispatched. \\
unpl.\ / und. & Among the loads a \emph{volumetric} plan dispatches,
  those \emph{proved} unplaceable and those left \emph{undecided}
  within the limit; an undecided load is counted on neither side. \\
$\Delta z$ & The cost of one method relative to another on the same
  instance, with the denominator named where it appears. \\
\bottomrule
\end{tabular}
\end{table}


\paragraph{Instances.}
Thirty-six instances form a full factorial over the products
$|N|\in\{4,8\}$, the horizon $|T|\in\{5,10\}$, the RTI types
$|R|\in\{2,3,4\}$ and the volumetric granularity $v\in\{60,90,120\}$,
named $I\_|N|\_|T|\_|R|\_v$, where $v$ is the number of erected RTIs
that fill one truck by volume, so larger $v$ means smaller RTIs and a
harder placement subproblem. Footprints are the five of the industrial
catalogue, from $800\times600$ to $2400\times1200$ mm, with heights
quantised to $50$ mm and bounded so that at least two layers fit under
the roof; from $|R|\ge3$ several types share a footprint, which lets a
stack mix types across layers. Demand deviates $\pm10\%$ from a mean of
roughly $2.5$ truckloads per period, and content weight by $\pm25\%$
between products sharing a type, which is what makes the crush limit
depend on the period's mix. The truck is a semi-trailer,
$13600\times2400\times3000$ mm, payload $24$ t, $W^{\max}=1200$ kg,
$\alpha=0.55$, $cK=1000$ per dispatch, $lt=1$, one period of coverage
stock. A size ladder follows the same rules over
$(|N|,|T|)\in\{(8,10),(12,12),(16,15),(24,20)\}$ crossed with the
truckloads consumed per period, $2.5$, $5$ or $8$, at $|R|=2$ and
$v=90$; its name's suffix is that figure times ten, not $K$.

\paragraph{Verification, and a second implementation.}
Every plan below is re-derived from its dispatch list alone and checked
truck by truck against \eqref{eq:col-geom}--\eqref{eq:col-agg} by a
routine independent of the search and of the oracle that built the
column, re-deriving each slot's weight from the fill the slot carries
and rejecting a fill that does not describe the stack it sits in. It
shares the full-truckload predicate with the oracle and the partition
floor, so it catches a defect in a solver's \emph{satisfaction} of that
definition and not one in the definition. The method has in addition
been re-implemented in a second language and solver stack and checked
layer by layer against the first, agreeing on $S\_2\_2\_2$ and
$I\_4\_5\_2\_60$ on every priced column, on every master dual to
$0.00\mathrm{e}{+}00$, on the certified bound to the digit and on the
whole search trajectory of the instance where the search exhausts ---
agreement that is partly by construction, a port having been tested
against the reference, and on $I\_4\_5\_2\_60$ the two plans in fact
differ.

Which of the two produced which table is stated here rather than left
to be inferred. Table~\ref{tab:volumetric} is the second
implementation, on the equal budgets its caption declares. Every other
table is the first. The two are not interchangeable at the sizes of the
size ladder: column generation there runs as the ladder of
Section~\ref{sec:pricing}, a heuristic rung ahead of the exact oracle,
and the second implementation has only the exact rung. Where the
heuristic rung contributes nothing --- the three smallest ladder rungs,
measured below --- the two agree or the second is slightly ahead; where
it contributes most of the trucks, the second implementation is behind
on one rung and returns no plan on the largest. A certificate is
unaffected either way, an exact sweep that returns nothing being a
certificate on its own, so this is a statement about which plans each
can find and not about what either may claim.

\subsection{What the volumetric side gives}
\label{sec:res-vol}

% Method rows: results/julia/parity_vol.csv, the volumetric parity campaign of
% 2026-10-27 -- the Julia implementation on all six instances at the compact
% reference's own budget (3600 s, 20 GB, 10 threads), commit 10e633c.  Every
% row carries root_proved=true, so none of them rests on the unproved-root
% bound defect fixed in e85c70c; that fix changes only the case where no
% certificate exists, and here one always does.  The t column is the `seconds`
% field, the method's own wall clock.  It lands at 3062-3133 s, not 3600:
% tree.jl:601-602 gives the search (1 - frontier_share) * time_limit = 3060 s
% and reserves the remainder for the frontier sweep, which returns in 2.3 to
% 72.5 s on these six.  The shortfall is the reserve going unspent, not a
% phase refused.
% Compact rows: results/compact_ref/compact_ref_<instance>.json, each with its
% own config block: 3600 s (1800 on I_8_5_2_90), MIP gap 0, 10 threads, 20 GB,
% seed 0.
% NOTE on the compact t column: every row prints the files' `wall` field
% (0.198167, 3268.92, 1349.52, 1121.83, 2093.27, 2403.21), not `runtime`
% (0.185, 3242.96, 1349.44, 1112.70, 2087.63, 2399.01).  One field for the
% whole column, and `wall` is the comparable one: the method's own time is
% wall-clock, so setting it against the compact model's solver time would
% flatter the method by the model-build time it does not charge itself.  The
% two fields differ by under 1% on every row, so nothing in the comparison
% turns on the choice -- but a column read off two fields is indefensible.
\begin{table}[!htbp]
\centering
\small
\caption{BP-3D-SCPTOP under its \emph{volumetric} oracle against the
compact SCPTOP formulation, both solving the \emph{volumetric} problem
of Appendix~\ref{appendix:A}, so that the two bracket one optimum.
\textbf{This table is a check on the procedure, not a contest}: the
volumetric problem is the one case here where an independent optimum
exists, and the question each row answers is whether the procedure
reaches it. $z$ is the incumbent at the stop; \textsc{opt} marks a
proved optimum, $^{\mathrm{m}}$ a stop on the memory cap,
$^{\mathrm{t}}$ a stop on time. ``bound'' is Gurobi's dual bound for
the compact model and the certified master-LP bound of
Section~\ref{sec:bounds} for the method. $\Delta z$ is
$(z_{\mathrm{DW}}-z_{\mathrm{compact}})/z_{\mathrm{compact}}$;
negative means the method's plan is cheaper. Both sides had $3\,600$ s,
$20$ GB and ten threads --- except that the compact model was given
$1\,800$ s on $I\_8\_5\_2\_90$, where it proved the optimum in $0.2$ s
and the budget never bound. The method's wall clock stops
short of $3\,600$ s by construction: the search is given
$(1-0.15)\times 3\,600 = 3\,060$ s and the rest is reserved for the
frontier sweep that bounds the regions left open, which returns here in
$2$ to $73$ s rather than spending its reserve. Reserved out of the
budget and not on top of it --- the recorded runs this table replaces
finished at $1\,817$ to $1\,897$ s against an $1\,800$ s budget,
and a row whose runtime exceeds the budget it is compared at is not a
comparison. No plan in this table is verified truck by truck: a volumetric
column carries no layout.}
\label{tab:volumetric}
\resizebox{\textwidth}{!}{%
\begin{tabular}{lrrrrlrrrrr}
\toprule
 & \multicolumn{4}{c}{compact SCPTOP} & \multicolumn{4}{c}{BP-3D-SCPTOP (volumetric oracle)} & & \\
\cmidrule(lr){2-5}\cmidrule(lr){6-9}
Instance & $z$ & bound & trucks & $t$ (s) & $z$ & bound (cert.) & trucks & $t$ (s) & $\Delta z$ (\%) & tighter bound \\
\midrule
$I\_8\_5\_2\_90$ & 18\,796.96\,\textsc{opt} & 18\,796.96 & 12 & 0.198 & 18\,796.96$^{\mathrm{t}}$ & 18\,795.56 & 12 & 3064 & $0.0000$ & 18\,796.96 (compact) \\
$I\_8\_10\_2\_90\_k25$ & 38\,855.65$^{\mathrm{m}}$ & 38\,851.93 & 26 & 3269 & 38\,853.17$^{\mathrm{t}}$ & 38\,791.19 & 26 & 3062 & $-0.0064$ & 38\,851.93 (compact) \\
$I\_16\_15\_2\_90\_k25$ & 68\,139.17\,\textsc{opt} & 68\,139.17 & 45 & 1350 & 68\,139.21$^{\mathrm{t}}$ & 68\,131.48 & 45 & 3117 & $+0.0001$ & 68\,139.17 (compact) \\
$I\_16\_15\_2\_90\_k50$ & 134\,346.91$^{\mathrm{m}}$ & 133\,417.47 & 89 & 1122 & 134\,507.47$^{\mathrm{t}}$ & 133\,484.56 & 89 & 3073 & $+0.1195$ & 133\,484.56 (DW) \\
$I\_24\_20\_2\_90\_k50$ & 200\,266.21$^{\mathrm{m}}$ & 199\,518.04 & 137 & 2093 & 200\,048.69$^{\mathrm{t}}$ & 200\,013.57 & 137 & 3071 & $-0.1086$ & 200\,013.57 (DW) \\
$I\_24\_20\_2\_90\_k80$ & 282\,101.77$^{\mathrm{m}}$ & 280\,823.53 & 189 & 2403 & 281\,233.28$^{\mathrm{t}}$ & 280\,910.22 & 188 & 3133 & $-0.3079$ & 280\,910.22 (DW) \\
\bottomrule
\end{tabular}}
\end{table}


The volumetric problem of Appendix~\ref{appendix:A} earns its place
in this study for one reason: it is the only part of the problem for
which an independent optimum exists, because a compact formulation of
it can be solved exactly on the smaller instances. What
Table~\ref{tab:volumetric} measures is therefore not which method
wins but whether the procedure arrives where it should on a case whose
answer is known. Both sides were given the same $3\,600$ s, $20$ GB
and ten threads.

On the two instances the compact model closes, it does. On
$I\_8\_5\_2\_90$ the procedure returns $18\,796.96$ with twelve
trucks --- the proved optimum, to the cent, reached from a different
formulation by a different search. On $I\_16\_15\_2\_90\_k25$ it
returns $68\,139.21$ against a proved $68\,139.17$, also at
forty-five trucks: four cents on sixty-eight thousand, $0.0001\%$
above the truth, with its own certified bound $0.0113\%$ below it. The
procedure brackets a known optimum from both sides to within about a
hundredth of a percent. That is the check this subsection exists to
report, and it passes.

Where the compact model works it is also far the faster, at $0.2$ s
and $1\,350$ s against the procedure's roughly $3\,070$ s. On a
volumetric problem small enough to write down whole, writing it down
whole is the better plan; the decomposition is built for the problem
the compact model cannot express, and Section~\ref{sec:res-3d} is
where that is asked of it.

On the four instances where the compact model exhausts its memory cap
with its dual bound short of its incumbent, the two agree closely
enough to say the procedure is solving the same problem well: costs
within $0.31\%$ on every row, the same fleet on three and one truck
fewer on $I\_24\_20\_2\_90\_k80$. The procedure's plan is cheaper
on three and dearer on one, by $0.1195\%$ on
$I\_16\_15\_2\_90\_k50$, the row where it opened twelve nodes in
the hour and left its own interval $0.76\%$ wide. At this accuracy the
two are measuring each other, and neither row is a ranking.

Where the decomposition does separate itself is the bound: on the
three largest instances its certificate is the tighter of the two, by
$67$, $496$ and $87$ cost units. The compact model's dual bound
degrades as its memory cap approaches, while the master LP over a
finite pool certifies whatever the pool supports.

One such statement needs no comparison of plans at all. On
$I\_12\_12\_2\_90\_k25$ the method certifies $52\,105.31$ --- a
full sweep in which every pair priced optimal and nothing priced out,
over the volumetric relaxation and so valid for the 3D problem by
Proposition~\ref{prop:inclusion} --- $274$ s into its run, where the
compact volumetric model still stood at a dual bound of $51\,944.22$
when its $600$ s expired, both at ten threads and a $20$ GB cap. Both
bound the same relaxation, and the decomposition's is $0.310\%$
tighter.

\subsection{What the 3D side gives}
\label{sec:res-3d}

% results/fleet_tree_3d/tree_3d_<instance>_table_<budget>.json, every row commit
% 73411b9, dirty false, 10 threads, 20 GB per model, seed 0, run sequentially;
% vol_lb_proved true and plan_verified true with no complaint on all five.
% Every \InFlight cell is replaced by the 3D ladder re-run in the Julia
% implementation at these same per-rung budgets (900/900/1800/1800/3600 s).
\begin{table}[!htbp]
\centering
\small
\caption{BP-3D-SCPTOP on the size ladder, every row on one commit
($73411\mathrm{b}9$), ten
threads, a $20$ GB per-model cap and seed $0$, run sequentially.
\emph{root} is the plan the construction returns before any node is
popped and \emph{tree} the improvement the search then makes on it,
both read from the same run's trace, so what branching contributes is
measured within a run rather than across commits. \emph{open} is the
regions still on the stack at the deadline and \emph{dropped} those
left with no bound covering them; both $0$ does \emph{not} mean the
space was refuted, a region closed with a certified bound below the
incumbent having been bounded and not excluded, which is why row 1
claims no optimum of its own. The bound is the certified volumetric
relaxation, valid for 3D by Proposition~\ref{prop:inclusion}; it is
loose wherever the crush limit binds, so the gap on the last two rows
is not an optimality gap. LB and \emph{gap} are the \emph{method's own}
certificate on every row, never a mixture: on rows 2 and 4 the compact
volumetric model supplies a tighter valid bound ($38\,851.93$ at a
memory stop, $68\,139.17$ proved optimal), which
Table~\ref{tab:volumetric} reports and Table~\ref{tab:3d-head} uses as
an interval end. Row 1 is the paper's only proved optimum and the proof
is not the method's: the compact volumetric relaxation of that instance
solves to $18\,796.96$ in $0.198$ s and this verified 3D plan attains
it. Row 1's instance is also row $I\_8\_5\_2\_90$ of
Table~\ref{tab:bench}: \textbf{two runs, on two commits} (there
$0\mathrm{c}0\mathrm{c}16\mathrm{c}$ at $900$ s), agreeing on $z$,
trucks and bound and differing at $68$ nodes against $67$ here, with the
gap printed there to two decimals rather than four. The twelve columns
are set as two blocks of the same five rows in the same order, so that
the table needs no reduction in type size.}
\label{tab:3d-onecommit}
\begin{tabular}{lrrlrr}
\toprule
Instance & $z$ & trucks & verified & LB (cert.) & gap \\
\midrule
$I\_8\_5\_2\_90$        & \InFlight{\textbf{18\,796.96}} & \InFlight{12}  & \InFlight{12 / 12}   & \InFlight{18\,795.54} & \InFlight{$0.0076\%$} \\
$I\_8\_10\_2\_90\_k25$  & \InFlight{38\,855.34} & \InFlight{26}  & \InFlight{26 / 26}   & \InFlight{38\,791.14} & \InFlight{$0.165\%$} \\
$I\_12\_12\_2\_90\_k25$ & \InFlight{52\,129.45} & \InFlight{34}  & \InFlight{34 / 34}   & \InFlight{52\,105.25} & \InFlight{$0.046\%$} \\
$I\_16\_15\_2\_90\_k25$ & \InFlight{78\,549.14} & \InFlight{54}  & \InFlight{54 / 54}   & \InFlight{68\,131.43} & \InFlight{$13.3\%$} \\
$I\_24\_20\_2\_90\_k50$ & \InFlight{230\,594.36} & \InFlight{160} & \InFlight{160 / 160} & \InFlight{200\,013.41} & \InFlight{$13.3\%$} \\
\bottomrule
\end{tabular}

\vspace{6pt}

\begin{tabular}{lrrrrrr}
\toprule
Instance & root ($t$) & tree & nodes & dropped & open & wall / budget (s) \\
\midrule
$I\_8\_5\_2\_90$        & \InFlight{18\,796.96 ($140$ s)} & \InFlight{$0.00\%$} & \InFlight{67} & \InFlight{0} & \InFlight{0}  & \InFlight{873 / 900} \\
$I\_8\_10\_2\_90\_k25$  & \InFlight{38\,865.30 ($136$ s)} & \InFlight{$0.03\%$} & \InFlight{68} & \InFlight{0} & \InFlight{35} & \InFlight{883 / 900} \\
$I\_12\_12\_2\_90\_k25$ & \InFlight{53\,195.84 ($459$ s)} & \InFlight{$\mathbf{2.00\%}$} & \InFlight{11} & \InFlight{0} & \InFlight{10} & \InFlight{1\,761 / 1\,800} \\
$I\_16\_15\_2\_90\_k25$ & \InFlight{78\,549.14 ($1\,181$ s)} & \InFlight{$0.00\%$} & \InFlight{1}  & \InFlight{0} & \InFlight{2}  & \InFlight{1\,746 / 1\,800} \\
$I\_24\_20\_2\_90\_k50$ & \InFlight{230\,594.36 ($1\,418$ s)} & \InFlight{$0.00\%$} & \InFlight{2} & \InFlight{0} & \InFlight{4} & \InFlight{3\,492 / 3\,600} \\
\bottomrule
\end{tabular}
\end{table}


% results/fleet_tree_3d/fleet_floor_<instance>_X<X>.json (commit dd5f039,
% region_infeasible true and region_proved true on all three rows).  The two
% I_16_15_2_90_k25 rows of the original table are dropped: they were certified
% against an earlier 53-truck plan and X<=53 was never tested, so they prove
% nothing about that instance.
\begin{table}[!htbp]
\centering
\small
\caption{Fleet floors certified offline. The volumetric relaxation's
master, restricted to $\sum_{g,t}n_{gt}\le X$, is priced to convergence
on one thread; \textsc{infeasible} means no plan with at most $X$ trucks
exists, in 3D as well, since every placeable load is volumetrically
feasible and the fleet row does not depend on the load model
(Proposition~\ref{prop:inclusion}). \emph{time} is the region
certification alone; the root certificate it rests on took $4$ to
$32$ s. These are the only proofs of minimality in this paper.}
\label{tab:fleetfloor}
\begin{tabular}{lrrlr}
\toprule
Instance & plan's trucks & $X$ tested & verdict & time (s) \\
\midrule
$I\_8\_5\_2\_90$        & 12 & $\le11$ & \textsc{infeasible}: 12 is the minimum fleet & 1.5 \\
$I\_8\_10\_2\_90\_k25$  & 26 & $\le25$ & \textsc{infeasible}: 26 is the minimum fleet & 3.9 \\
$I\_12\_12\_2\_90\_k25$ & 34 & $\le33$ & \textsc{infeasible}: 34 is the minimum fleet & 16.3 \\
\bottomrule
\end{tabular}
\end{table}


% Aggregated by hand from the same 36 files as Table \ref{tab:bench} and from
% results/loads_recheck.csv.  Totals verified by recomputation: 31/36 plans,
% 30 gap rows, mean 8.799%; 724 loads = 501 unplaceable + 174 undecided + 49
% placeable.
\begin{table}[!htbp]
\centering
\small
\caption{Table~\ref{tab:bench} by instance factor. \emph{plans} counts
the instances on which BP-3D-SCPTOP returned a verified plan inside the
$900$ s budget; the mean gap is over those instances and over the rows
whose bound was proved, and the count of each is given so that the
denominator is never in doubt. The last three columns are the
volumetric benchmark's dispatched loads: the share a constraint
programme \emph{proves} unplaceable, the share it leaves
\emph{undecided}, and the share it places. The proved share is a
\emph{floor}, since more oracle effort can only move loads out of the
undecided column.}
\label{tab:summary}
\begin{tabular}{llrrrrrr}
\toprule
factor & value & plans & mean gap (rows) & loads & unpl. & und. & placed \\
\midrule
$|N|$ & 4 & 16 / 18 & $8.9\%$ (16) & 363 & $67.2\%$ & $26.4\%$ & $6.3\%$ \\
 & 8 & 15 / 18 & $8.6\%$ (14) & 361 & $71.2\%$ & $21.6\%$ & $7.2\%$ \\
\midrule
$|T|$ & 5 & 17 / 18 & $8.4\%$ (17) & 219 & $56.2\%$ & $33.8\%$ & $10.0\%$ \\
 & 10 & 14 / 18 & $9.4\%$ (13) & 505 & $74.9\%$ & $19.8\%$ & $5.3\%$ \\
\midrule
$|R|$ & 2 & 10 / 12 & $9.7\%$ (10) & 246 & $79.3\%$ & $15.4\%$ & $5.3\%$ \\
 & 3 & 11 / 12 & $8.4\%$ (10) & 235 & $61.7\%$ & $29.4\%$ & $8.9\%$ \\
 & 4 & 10 / 12 & $8.3\%$ (10) & 243 & $66.3\%$ & $27.6\%$ & $6.2\%$ \\
\midrule
$v$ & 60 & 7 / 12 & $19.3\%$ (6) & 182 & $91.8\%$ & $2.7\%$ & $5.5\%$ \\
 & 90 & 12 / 12 & $5.0\%$ (12) & 239 & $55.2\%$ & $36.4\%$ & $8.4\%$ \\
 & 120 & 12 / 12 & $7.4\%$ (12) & 303 & $66.7\%$ & $27.1\%$ & $6.3\%$ \\
\midrule
all & & 31 / 36 & $8.8\%$ (30) & 724 & $\mathbf{69.2\%}$ & $24.0\%$ & $6.8\%$ \\
\bottomrule
\end{tabular}
\end{table}


On the thirty-six generated instances at a uniform $900$ s the method
returns a verified plan on thirty-one and a proved bound on
thirty-five, at a mean certified gap of $8.8\%$ over the thirty rows
with both. The five failures are not spread across the factorial:
all five have $v=60$, where one container is the largest fraction of a
trailer and a coarse integer loading structure leaves the construction
least room to cover a delivery vector with full trucks --- which is
also why a volumetric capacity is least adequate there
(Section~\ref{sec:res-abstraction}). On those five rows
\emph{no plan} is returned rather than an invalid one, which follows
from the formulation: a column the verifier would reject never enters
the pool.

The ladder takes the same method to twenty periods and
$\InFlight{160}$ trucks under the same guarantee. On the first three
rungs the truck count is in addition the minimum possible, the region
with one truck fewer being certified infeasible in $1.5$, $3.9$ and
$16.3$ s (Table~\ref{tab:fleetfloor}), so
what remains of those gaps, $\InFlight{0.0076\%}$--$\InFlight{0.165\%}$,
is inventory cost. The jump to
$\InFlight{13.3\%}$ on the last two rungs is attributed to the bound
and not to the plan, on the mechanistic evidence of
Section~\ref{sec:res-abstraction}; nothing here proves that
attribution, and those two plans are reported as verified upper bounds
under a stated budget. The claim about the gaps is therefore
conditional: they are small where the volumetric relaxation is nearly
exact, which is checkable in advance.

What branching contributes is uneven and is not averaged away. Read
within one run (Table~\ref{tab:3d-onecommit}), the tree improves the
root construction by $\InFlight{2.00\%}$ on the third rung, five node
improvements taking $\InFlight{53\,195.84}$ to $\InFlight{52\,129.45}$,
and by $\InFlight{0.03\%}$ or nothing on the other four: the pattern
follows what the root leaves the tree, and where the relaxation is
nearly exact the construction is already there. That the two generators
are complementary is measured under one dual trajectory --- the exact
oracle's columns carry $4.72$ distinct commodities on average against
the heuristic rung's $10.17$, which is Figure~\ref{fig:hull} in numbers
--- and in production that rung contributes $0$ trucks to each of the
three plans whose gap is under $0.2\%$, $0$ of $\InFlight{12}$, of
$\InFlight{26}$ and of $\InFlight{34}$, against $\InFlight{44}$ of
$\InFlight{54}$ and $\InFlight{122}$ of $\InFlight{160}$ on the two
where the crush limit binds.

% Phase clock: results/fleet_tree_3d/tree_3d_I_8_5_2_90_table_900.json,
% t_cg 864.16 of seconds 873.10.
% Packing census, summed over the 36 results/bench_tree/tree_3d_*.json:
% stats.pack_ok_shelf 24686, stats.pack_ok_cpsat 246,
% stats.distinct_parts 11866, stats.distinct_parts_unplaceable 4327,
% stats.distinct_parts_undecided 397.
The bound is cheap, which makes the division of labour of
Section~\ref{sec:bounds} affordable: over the forty-one runs reported
here the volumetric sweep returns in $2.5$ to $111.8$ s with every pair
proved on forty, the exception being the one row of
Table~\ref{tab:bench} that prints \emph{unproved} and is excluded from
every mean. The geometric sweep certified no region on any run. Where
the time goes depends on whether the crush limit binds: column
generation takes $\InFlight{864}$ s of the $\InFlight{873}$ s first
rung with the geometry free, while on the two crush-binding rungs the
partition and its placeability calls dominate. Over the thirty-six
instances the shelf design placed $24\,686$ loads and the constraint
programme $246$, which proved $4\,327$ of the $11\,866$ distinct
compositions proposed unplaceable and left $397$ undecided.

\subsection{The compact model}
\label{sec:res-compact}

% Compact 3D columns: results/compact_3d/compact3d_*.json (1800 s, MIP gap 0,
% 10 threads, 20 GB, seed 0, n_lc 3); row 1 is compact3d_rerun_I_8_5_2_90.json,
% taken after the layout-extraction fix, with the same objective, bound and stop
% as the first run but twelve loads that pass the verifier.  No compact 3D run
% exists for I_24_20_2_90_k50.  Method columns are Table \ref{tab:3d-onecommit}.
% Only the ends that a campaign replaces are marked: a compact volumetric dual
% bound (rows 2 and 4's lower end) is not in flight.
% The interval's lower end is the tighter bound THAT BOUNDS THE PROBLEM: the
% compact volumetric dual bound where one exists, else the method's own
% certificate -- never the compact 3D model's dual bound.
\begin{table}[!htbp]
\centering
\small
\caption{The compact 3D-SCPTOP model (a catalogue of
$n_{\mathrm{LC}}=3$ layouts shared by the trucks of a direction,
$1\,800$ s) against BP-3D-SCPTOP on the size ladder, the method's rows
being those of Table~\ref{tab:3d-onecommit}. \textbf{The budgets differ
on three rows}: the compact model had $1\,800$ s throughout, the method
$900$ s on the first two rows, $1\,800$ s on the next two --- which are
therefore the equal-budget pairs --- and $3\,600$ s on the last.
\emph{none} means no feasible solution was found inside the budget, and
all three such runs stopped on \emph{time} with peak resident sets of
$6.1$, $2.7$ and $4.5$ GB against the declared $20$ GB cap. The last
column is the interval in which $z^{\star}$ lies, taking the method's
plan as the upper end and, as the lower end, the tighter bound
\emph{that bounds the problem}: $n_{\mathrm{LC}}=3$ is an a-priori
restriction of Section~\ref{sec:model}, so the compact 3D model's dual
bound bounds that restriction and not $z^{\star}$ and appears in its
own column only, while a compact \emph{volumetric} dual bound is valid
for 3D by inclusion and is used instead where one exists. Row 1 is the
one optimum (Table~\ref{tab:3d-onecommit}).}
\label{tab:3d-head}
\resizebox{\textwidth}{!}{%
\begin{tabular}{lrrrlrrrrrl}
\toprule
& \multicolumn{4}{c}{compact 3D-SCPTOP, $1\,800$ s} & \multicolumn{5}{c}{BP-3D-SCPTOP} & \\
\cmidrule(lr){2-5}\cmidrule(lr){6-10}
Instance & $z$ & trucks & LB & stop & $z$ & trucks & LB (cert.) & gap & budget (s) & $z^{\star}$ in \\
\midrule
$I\_8\_5\_2\_90$ & 18\,816.13 & 12 & 18\,795.76 & \textsc{time} & \InFlight{18\,796.96} & \InFlight{12} & \InFlight{18\,795.54} & \InFlight{$0.0076\%$} & 900 & $[\InFlight{18\,796.96},\,\InFlight{18\,796.96}]$ \\
$I\_8\_10\_2\_90\_k25$ & \emph{none} & --- & 36\,458.75 & \textsc{time} & \InFlight{38\,855.34} & \InFlight{26} & \InFlight{38\,791.14} & \InFlight{$0.165\%$} & 900 & $[38\,851.93,\,\InFlight{38\,855.34}]$ \\
$I\_12\_12\_2\_90\_k25$ & \emph{none} & --- & 48\,516.86 & \textsc{time} & \InFlight{52\,129.45} & \InFlight{34} & \InFlight{52\,105.25} & \InFlight{$0.046\%$} & 1\,800 & $[\InFlight{52\,105.25},\,\InFlight{52\,129.45}]$ \\
$I\_16\_15\_2\_90\_k25$ & \emph{none} & --- & 64\,053.23 & \textsc{time} & \InFlight{78\,549.14} & \InFlight{54} & \InFlight{68\,131.43} & \InFlight{$13.3\%$} & 1\,800 & $[68\,139.17,\,\InFlight{78\,549.14}]$ \\
$I\_24\_20\_2\_90\_k50$ & \multicolumn{4}{c}{not run} & \InFlight{230\,594.36} & \InFlight{160} & \InFlight{200\,013.41} & \InFlight{$13.3\%$} & 3\,600 & $[\InFlight{200\,013.41},\,\InFlight{230\,594.36}]$ \\
\bottomrule
\end{tabular}}
\end{table}


A compact formulation of 3D-SCPTOP cannot serve as a yardstick beyond
five periods. Written with a catalogue of $n_{\mathrm{LC}}=3$ layouts
shared by the trucks of a direction --- a \emph{restriction} of
Section~\ref{sec:model} chosen a priori, so its dual bound bounds the
restriction and not $z^{\star}$ --- it returns a plan on one of four
ladder instances. On $I\_8\_5\_2\_90$, at $67\,674$ variables and
$92\,092$ constraints, it stops on its $1\,800$ s limit with no phantom
freight and twelve loads that pass the independent verifier,
$\InFlight{0.102\%}$ above the method's plan at twice the method's
budget. From ten periods on, at $119\,424$ to $322\,734$ variables, it
% Peak resident sets: results/compact_3d/compact3d_I_8_10_2_90_k25.json
% 6.1205, compact3d_I_12_12_2_90_k25.json 2.6528,
% compact3d_I_16_15_2_90_k25.json 4.4501 (field peak_rss_gb), in the row
% order of Table \ref{tab:3d-head}.
returns \emph{no feasible solution at all} in $1\,800$ s, all three
stopping on time and not on memory --- peak resident sets of $6.1$,
$2.7$ and $4.5$ GB against a declared $20$ GB cap. Four rungs are four
rungs; beyond five periods the compact model supplies neither a plan
nor a bound that bounds the problem.

$I\_8\_5\_2\_90$ is also the one instance on which this paper reports
an optimum, and the proof is not the method's. The compact
\emph{volumetric} relaxation solves it to proved optimality at
$18\,796.96$ in $0.198$ s, every placeable load is volumetrically
feasible, and the verified 3D plan attains that value; the method's own
% The three leaf certifications are a separate run:
% results/fleet_tree_3d/tree_3d_I_8_5_2_90_leafcert_900.json, leaves_certified 3,
% log lines 151-153 "leaf certification 18795.57/18795.54/18795.54
% proved=False; not refuted".  Table \ref{tab:3d-onecommit}'s row 1 is
% tree_3d_I_8_5_2_90_table_900.json, whose bound is the volumetric certificate.
attempt failed: geometric certification of the three regions its search
could not refute returned $18\,795.57$, $18\,795.54$ and
$18\,795.54$, none of the three \emph{proved}. The
proof is available because this instance's two container families have
coinciding height- and crush-limited layer counts --- exactly the
condition under which Section~\ref{sec:res-abstraction} finds a
volumetric capacity adequate.

\subsection{Is a volumetric capacity adequate?}
\label{sec:res-abstraction}

% results/volumetric_loads_recheck.json (pack_tl 60 s, oracle pack_witness, every
% placed witness re-verified by verify_column): I_16_15_2_90_k25 dispatched 45,
% unplaceable 41, undecided 0, source compact_ref OPTIMAL 68139.17;
% I_24_20_2_90_k50 dispatched 137, unplaceable 119, undecided 0, source
% compact_ref MEM_LIMIT 200266.21.
% results/repacked_volumetric.json for rows 2-4 (loadlocal_trucks 86 / 256,
% repacked_trucks 74 / 212, method_trucks 54 / 160; pairs_unpacked 0 and
% complaints 0 on both; over_delivered_units 80 / 137).
% The method rows are replaced by the 3D ladder re-run (900-3600 s per rung).
\begin{table}[!t]
\centering
\small
\caption{What volumetric planning costs, constructed rather than
estimated: trucks dispatched on the two instances where the per-stack
crush limit binds. Row 1 is the volumetric benchmark's own plan, which
is not executable --- $41$ of its $45$ dispatched loads on the first
instance and $119$ of its $137$ on the second are \emph{proved}
unplaceable with \emph{none} left undecided, so those shares carry no
undecided caveat. Row 2 splits each unplaceable load where it stands,
with no consolidation; it is reported for contrast and is \emph{not}
the comparison, consolidation being worth twelve and forty-four trucks.
Row 3 is the comparison: the volumetric plan's delivery quantities
kept, each period re-packed with the same partition-and-pack machinery
the method uses, every resulting truck verified, and at least the
original quantities delivered --- by $80$ and $137$ units, the
partition being a cover. Row 4 is BP-3D-SCPTOP, which also chooses the
quantities. The volumetric plan on $I\_24\_20\_2\_90\_k50$ is a
\emph{memory stop} and not an optimum, so a better one might repair into
fewer than $212$ trucks and that column's margin is overstated by an
unquantified amount in the method's favour; the
$I\_16\_15\_2\_90\_k25$ plan is the proved volumetric optimum and
carries no such caveat.}
\label{tab:repack}
\begin{tabular}{lrr}
\toprule
 & $I\_16\_15\_2\_90\_k25$ & $I\_24\_20\_2\_90\_k50$ \\
\midrule
volumetric plan (not executable)      & 45  & 137 \\
repaired load by load (contrast only) & 86  & 256 \\
\textbf{re-packed, quantities kept}   & \textbf{74} & \textbf{212} \\
\textbf{BP-3D-SCPTOP} (which also chooses the quantities) & \textbf{\InFlight{54}} & \textbf{\InFlight{160}} \\
\midrule
saving on trucks / on cost & \InFlight{$27.0\%$ / $19.1\%$} & \InFlight{$24.5\%$ / $16.2\%$} \\
\bottomrule
\end{tabular}
\end{table}


The question is answered by proof: every load the volumetric benchmark
dispatches is submitted to an oracle asking whether that exact set of
RTIs can be stacked and placed in a trailer at all, with the true slot
bound \eqref{sb:eq:Ibound} in force so that an infeasible verdict is a
proof. Of the $724$ loads the benchmark dispatches over the thirty-six
generated instances, $501$ --- $69.2\%$ --- are \emph{proved}
unplaceable. That figure is a floor and it is the one to read: $174$
($24.0\%$) were left \emph{undecided} within the limit and only $49$
($6.8\%$) proved placeable, so more oracle effort can only move loads
out of the undecided column, and the net of the two pieces of evidence
below is unknown. Re-running the same oracle at two and a half times the limit
decided twenty more loads, all twenty placeable, leaving the proved
count at $501$; and the oracle carries no payload row, so an
over-payload load comes back undecided rather than refuted. The share
of \emph{decided} loads that cannot be placed is correspondingly an
upper reading, drifting downward with compute --- $94.5\%$ at the
shorter limit, $91.1\%$ at the longer --- and rising to $91.8\%$ at
$v=60$ where only $2.7\%$ are undecided (Table~\ref{tab:summary}).

The two ladder instances where the crush limit binds carry no undecided
caveat at all: on $I\_16\_15\_2\_90\_k25$, whose volumetric plan is the
\emph{proved} optimum, $41$ of its $45$ dispatched loads are proved
unplaceable and four placeable, and on $I\_24\_20\_2\_90\_k50$ it is
$119$ of $137$. Those two went to the stronger oracle of
Section~\ref{sec:geo3d} --- a constructive witness first, the
constraint programme only for refutations --- which left nothing
undecided; the thirty-six went to the constraint programme alone, which
is why a quarter of their loads are. A volumetric optimum can therefore
be a plan nine tenths of which cannot be loaded, and the error has a
direction: a volumetric capacity ignores the remainders that footprints
and layer heights leave unusable.

What that costs is constructed rather than estimated
(Table~\ref{tab:repack}): each volumetric plan's delivery quantities
are kept and each period's freight re-packed with the same
partition-and-pack machinery the method uses, every resulting truck
verified. The re-packed plans need $74$ and $212$ trucks where the
method plans $\InFlight{54}$ and $\InFlight{160}$ directly --- savings
of $\InFlight{27.0\%}$ and $\InFlight{24.5\%}$ on trucks and
$\InFlight{19.1\%}$ and $\InFlight{16.2\%}$ on cost, each denominator
the re-packed plan's own figure. Two caveats travel with them: the
re-packed plan delivers \emph{at least} the volumetric quantities, by
$80$ and $137$ units, the partition being a cover; and the
$I\_24\_20\_2\_90\_k50$ volumetric plan is a memory stop rather than an
optimum, so a better one might repair into fewer than $212$ trucks and
that row's margin is overstated in the method's favour. The
$I\_16\_15\_2\_90\_k25$ row, the proved optimum, carries no such
caveat.

Deflating the capacity is the obvious remedy, and one step is measured:
adding the stack-count rows \eqref{eq:stackrows} to the compact
volumetric model, valid for every placeable load, raises its answer on
$I\_16\_15\_2\_90\_k25$ from forty-five trucks to forty-eight and its
dual bound to $70\,435.04$ after $600$ s --- $2\,296$ above that
model's own bound of $68\,139.17$ --- while changing nothing where the
crush limit does not bind. Whether those forty-eight loads can be
placed has not been tested, the stack rows being necessary for
placeability and not sufficient, so the stack-aware model is not yet a
baseline.

\paragraph{The mechanism, and a test a planner can run.}
All of it has one cause, measured directly on
$I\_16\_15\_2\_90\_k25$: the payload admits $57$ full RTIs of product
$P_{1}$ in one truck while $32$ is the largest number that can be
placed, a proved bound from pricing that commodity alone to optimality.
The container is $650$ mm tall, so four layers fit under the
$3\,000$ mm roof, while $P_{1}$ at $419$ kg full fits only two under
the $1\,200$ kg crush limit, as do six of the family's eight products.
A volumetric fill therefore credits a truck with roughly twice that
family's placeable count, which is also why the certified bound on the
two crush-binding rungs sits $\InFlight{13.3\%}$ below the best
verified plan.

That mechanism can be checked before anything is solved, since it needs
only the catalogue and the trailer. For one container type, let the
\emph{analytic credit ratio} $\kappa$ be the count a volumetric
capacity credits over an \emph{analytic} count of how many can be
carried,
\begin{equation}
\kappa_{r}=
\frac{\min\{\Gamma_{K}/\gamma_{pr},\;pK^{w}/\bar w_{r}\}}
     {\min\{s_{r}\min(L^{h}_{r},L^{c}_{r}),\;pK^{w}/\bar w_{r}\}},
\qquad
L^{h}_{r}=\Bigl\lfloor\tfrac{rK}{r_{pr}}\Bigr\rfloor,
\quad
L^{c}_{r}=\Bigl\lfloor\tfrac{W^{\max}}{\bar w_{r}}\Bigr\rfloor,
\label{eq:kappa}
\end{equation}
with $\bar w_{r}$ the catalogued gross weight of one full RTI and
$s_{r}$ the number of its footprints a single orientation tiles onto
the floor. The denominator is a grid count with no packing oracle run,
hence an \emph{upper} estimate of what fits, so $\kappa>1$ is a
conservative diagnosis of over-credit. Over the fifteen types of the
industrial catalogue $\kappa$ has median $1.000$, ranges to $2.000$ and
exceeds $1.20$ on four types, and \emph{no} type has $\kappa<1$: the
error cannot run in the planner's favour, and a plan over a mixed
catalogue fails on its worst container. The census instances are
\emph{generated}, with $\kappa$ in a range the generator controls, so
this is a realism check on that range rather than a claim that the
census shares transfer to a particular plant.